% CREATED BY MAGNUS GUSTAVER, 2020
\chapter{Results}
\label{chap:results}

This chapter evaluates NanoMem as a memory evidence provider. The experiments
ask whether the system can construct compact, source-grounded evidence for
long-term memory questions, whether the benefit is concentrated in temporal
and multi-session settings, and which parts of the method explain the observed
behavior. The reported output of NanoMem is the evidence passed to a frozen
answer model, not the final answer generated by the memory module itself.

The chapter is organized as follows. Section~\ref{sec:results-protocol}
summarizes the evaluation protocol, datasets, baselines, and metrics.
Section~\ref{sec:results-main} reports the main benchmark results.
Section~\ref{sec:results-compactness} analyzes compactness. Sections
\ref{sec:results-architecture} and~\ref{sec:results-reward} report component
and reward ablations. Section~\ref{sec:results-cases} gives qualitative
traces and Section~\ref{sec:results-errors} summarizes the observed failure
types.

\begin{figure}[H]
  \centering
  \fbox{\begin{minipage}{0.92\textwidth}
    \centering
    \vspace{0.8em}
    \textcolor{red}{TODO: Replace this placeholder with the planned Results
    experiment map. It should connect RQ3 to main benchmark results,
    compactness, architecture ablations, reward diagnostics, case studies, and
    error analysis.}
    \vspace{0.8em}
  \end{minipage}}
  \caption{Planned map of the Results chapter. Each experiment group targets a
  different claim about NanoMem: benchmark accuracy, compact evidence,
  component necessity, learned stop/continue behavior, and failure boundaries.}
  \label{fig:results-map}
\end{figure}

\section{Evaluation Protocol}
\label{sec:results-protocol}

The evaluation uses three complementary benchmark settings. LoCoMo evaluates
long-term multi-session dialogue memory with single-hop, multi-hop, temporal,
and open-domain questions \cite{maharana2024locomo}. LongMemEval evaluates
long-term assistant memory over noisier histories, including single-session
questions, knowledge updates, temporal reasoning, and multi-session reasoning
\cite{wu2024longmemeval}. ChainMem is the diagnostic benchmark constructed in
this project for hidden-dependency temporal bridge questions: the query cannot
directly name the destination evidence, so the system must first recover a
source event and then use the discovered time, entity, or condition to search
again.

The common evaluation boundary is fixed across NanoMem variants. NanoMem
returns structured evidence and a verdict. A downstream Qwen3.5-9B answerer
then generates the final answer from the query and evidence, and a GPT-5.1
judge compares the prediction against the gold answer. This boundary matters:
an answer can fail because retrieval missed the source, because the Synthesizer
emitted incomplete evidence, because the verdict stopped too early, or because
the downstream answerer made an error despite correct evidence. The metrics
below therefore report both answer accuracy and evidence-side behavior.

\begin{table}[H]
  \centering
  \caption{Datasets used in the Results chapter.}
  \label{tab:results-datasets}
  \scriptsize
  \begin{tabular}{p{0.18\textwidth}p{0.18\textwidth}p{0.18\textwidth}p{0.14\textwidth}p{0.20\textwidth}}
    \hline\hline
    Benchmark & Main question types & Temporal grounding & Hidden dependency & Role in this chapter \\
    \hline
    LoCoMo & Single-hop, multi-hop, temporal, open-domain & Yes & Sometimes & Standard long-dialogue memory evaluation \\
    LongMemEval & SSU, SSA, KU, TR, MR & Yes & Sometimes & Noisier long assistant-memory evaluation \\
    ChainMem & Normal and hard temporal bridge questions & Yes & Yes & Diagnostic test for evidence addressability gap \\
    \hline\hline
  \end{tabular}
\end{table}

\begin{table}[H]
  \centering
  \caption{Training and held-out evaluation split used for NanoMem.}
  \label{tab:results-splits}
  \scriptsize
  \begin{tabular}{p{0.22\textwidth}rrrrp{0.20\textwidth}}
    \hline\hline
    Data source & Raw & Filtered & Train & Held-out & Leakage control \\
    \hline
    LoCoMo & 1986 & 1536 & 401 & 655 & Conversation-level split \\
    LongMemEval & 500 & 469 & 0 & 469 & Held out from GRPO training \\
    ChainMem & 500 & 374 & 99 & 275 & Filtered held-out diagnostic examples \\
    \hline\hline
  \end{tabular}
\end{table}

Baselines are grouped by the stage at which they operate. Full-context
baselines give the complete history directly to GPT-5.1 or Qwen3.5-9B and
represent strong long-context references. Write-time memory baselines such as
Mem0 and A-Mem store extracted or organized memory records before the future
query is known \cite{chhikara2025mem0,xu2025amem}. Search-time baselines such
as MemR$^3$ and Search-R1 adapt retrieval after the query arrives
\cite{du2025memr3,jin2025searchr1}. Memory-T1 is included as a temporal memory
baseline because it explicitly targets temporal evidence selection
\cite{du2026memoryt1}. NanoMem is reported in two versions: NanoMem-Vanilla,
which uses the same framework without the GRPO-trained Synthesizer, and
NanoMem-GRPO, which uses the trained Synthesizer.

\begin{table}[H]
  \centering
  \caption{Metrics reported in this chapter.}
  \label{tab:results-metrics}
  \scriptsize
  \begin{tabular}{p{0.20\textwidth}p{0.25\textwidth}p{0.25\textwidth}p{0.18\textwidth}}
    \hline\hline
    Metric & Computed from & Interpretation & Main caveat \\
    \hline
    Accuracy & GPT-5.1 judge over predicted and gold answers & Downstream task success & Depends on judge and answer model \\
    Output tokens & Evidence or returned memory context & Compactness of information passed downstream & Not identical to full input cost for all baselines \\
    Retrieval rounds & NanoMem trace & Amount of iterative search used & Fewer rounds are not always better \\
    Verdict accuracy & Sufficiency decisions against evidence coverage diagnostics & Stop/continue quality & Requires reliable gold evidence coverage \\
    \hline\hline
  \end{tabular}
\end{table}

\section{Main Benchmark Results}
\label{sec:results-main}

Table~\ref{tab:results-main-summary} gives a compact view of the main
benchmark outcomes. The most important pattern is not that NanoMem-GRPO wins
every column. Rather, its gains are concentrated where the thesis predicts
them: multi-hop, temporal reasoning, and hidden-dependency bridge questions.
These are the cases in which a one-shot query often lacks the actual retrieval
address.

\begin{table}[H]
  \centering
  \caption{Main benchmark summary. Accuracy is reported as percentage points.}
  \label{tab:results-main-summary}
  \scriptsize
  \resizebox{\textwidth}{!}{%
  \begin{tabular}{p{0.23\textwidth}rrr|rrr|rrr}
    \hline\hline
    Method & L-MH & L-Temp & L-All & LME-TR & LME-MR & LME-All & CM-N & CM-H & CM-All \\
    \hline
    GPT-5.1 full context & 88.7 & 76.0 & 87.7 & 39.8 & 44.4 & 63.6 & 25.0 & 7.0 & 16.0 \\
    Qwen3.5-9B full context & 84.0 & 55.5 & 80.8 & 39.8 & 45.1 & 63.6 & 25.0 & 16.0 & 20.5 \\
    Mem0 & 79.3 & 43.0 & 71.0 & 54.2 & 46.2 & 59.2 & 2.0 & 1.0 & 1.5 \\
    NanoMem-Vanilla & 84.0 & 76.6 & 85.0 & 75.8 & 70.7 & 78.7 & 47.0 & 39.0 & 43.0 \\
    NanoMem-GRPO & \textbf{89.0} & \textbf{87.5} & \textbf{88.6} & \textbf{86.4} & \textbf{71.4} & \textbf{83.8} & \textbf{57.0} & \textbf{47.0} & \textbf{51.5} \\
    \hline\hline
  \end{tabular}
  }
\end{table}

On LoCoMo, NanoMem-GRPO reaches 88.6\% overall accuracy, slightly above the
GPT-5.1 full-context reference at 87.7\%. The more informative comparison is
category-level. NanoMem-GRPO improves over NanoMem-Vanilla on multi-hop
questions from 84.0\% to 89.0\%, and on temporal questions from 76.6\% to
87.5\%. This supports the claim that the Planner--Retriever--Synthesizer loop
is most useful when the answer depends on a relation across events or times,
rather than on a single directly named fact.

\begin{table}[H]
  \centering
  \caption{LoCoMo category-level results. Accuracy is percentage; Tok. is mean
  output tokens for the returned evidence or memory context.}
  \label{tab:results-locomo}
  \scriptsize
  \resizebox{\textwidth}{!}{%
  \begin{tabular}{p{0.23\textwidth}rrrrrrrrrr}
    \hline\hline
    Method & SH & Tok. & MH & Tok. & Temp. & Tok. & Open & Tok. & All & Tok. \\
    \hline
    GPT-5.1 full context & 94.2 & 25 & 88.7 & 32 & 76.0 & 27 & 67.7 & 36 & 87.7 & 27 \\
    Qwen3.5-9B full context & 92.2 & 18 & 84.0 & 24 & 55.5 & 16 & 56.2 & 39 & 80.8 & 20 \\
    Mem0 & 79.5 & 675 & 79.3 & 650 & 43.0 & 650 & 67.1 & 675 & 71.0 & 650 \\
    NanoMem-Vanilla & 90.5 & 234 & 84.0 & 364 & 76.6 & 225 & 67.4 & 375 & 85.0 & 264 \\
    NanoMem-GRPO & \textbf{91.4} & 99 & \textbf{89.0} & 138 & \textbf{87.5} & 88 & 65.2 & 152 & \textbf{88.6} & 107 \\
    \hline\hline
  \end{tabular}
  }
\end{table}

On LongMemEval, the difference between full-context answering and evidence
selection is larger. NanoMem-GRPO reaches 83.8\% overall accuracy, compared
with 63.6\% for GPT-5.1 full context. This result should not be read as a
general claim that the full-context model is weak. It shows that in long,
noisy, multi-session histories, selecting and synthesizing relevant evidence
can be more effective than asking the answer model to process the whole
history directly. The temporal reasoning category is particularly strong:
NanoMem-GRPO reaches 86.4\%, while NanoMem-Vanilla reaches 75.8\%.

\begin{table}[H]
  \centering
  \caption{LongMemEval category-level results. SSU is single-session user, SSA
  is single-session assistant, KU is knowledge update, TR is temporal reasoning,
  and MR is multi-session reasoning.}
  \label{tab:results-longmemeval}
  \scriptsize
  \resizebox{\textwidth}{!}{%
  \begin{tabular}{p{0.20\textwidth}rrrrrrrrrrrr}
    \hline\hline
    Method & SSU & Tok. & SSA & Tok. & KU & Tok. & TR & Tok. & MR & Tok. & All & Tok. \\
    \hline
    GPT-5.1 full context & 95.7 & 31 & 96.4 & 48 & 84.6 & 29 & 39.8 & 52 & 44.4 & 52 & 63.6 & 42 \\
    Qwen3.5-9B full context & 97.1 & 76 & 100.0 & 74 & 79.5 & 148 & 39.8 & 255 & 45.1 & 189 & 63.6 & 172 \\
    NanoMem-Vanilla & 94.3 & 104 & 94.6 & 245 & 71.8 & 357 & 75.8 & 417 & 70.7 & 462 & 78.7 & 353 \\
    NanoMem-GRPO & 95.7 & 81 & 98.2 & 94 & 79.5 & 109 & \textbf{86.4} & 124 & \textbf{71.4} & 149 & \textbf{83.8} & 119 \\
    \hline\hline
  \end{tabular}
  }
\end{table}

ChainMem is harder in absolute terms. NanoMem-GRPO reaches 51.5\% overall
accuracy, compared with 43.0\% for NanoMem-Vanilla and far below saturation.
This lower ceiling is expected: the benchmark is designed so that the answer
session is hidden behind a bridge event, condition, or time. The gain therefore
matters more than the absolute number. NanoMem's iterative evidence state helps
because the first round can expose a new retrieval cue that was absent from the
original query.

\begin{table}[H]
  \centering
  \caption{ChainMem results by difficulty and terminal type. Normal and Hard
  split the bridge difficulty; entity, condition, and time describe the final
  answer type.}
  \label{tab:results-chainmem}
  \scriptsize
  \resizebox{\textwidth}{!}{%
  \begin{tabular}{p{0.22\textwidth}rrrrrrrrrrrrrr}
    \hline\hline
    Method & Norm. & Tok. & Hard & Tok. & Ent. & Tok. & Cond. & Tok. & Time & Tok. & All & Tok. \\
    \hline
    NanoMem-Vanilla & 42.5 & 318 & 53.3 & 312 & 46.7 & 267 & 40.0 & 258 & 33.3 & 239 & 43.0 & 287 \\
    NanoMem-GRPO & \textbf{52.5} & 221 & \textbf{66.7} & 266 & \textbf{53.3} & 226 & \textbf{47.5} & 216 & \textbf{46.7} & 256 & \textbf{51.5} & 234 \\
    \hline\hline
  \end{tabular}
  }
\end{table}

\begin{figure}[H]
  \centering
  \fbox{\begin{minipage}{0.92\textwidth}
    \centering
    \vspace{0.8em}
    \textcolor{red}{TODO: Replace this placeholder with a cross-benchmark
    accuracy--compactness plot. Plot NanoMem-GRPO, NanoMem-Vanilla, full
    context, and representative memory baselines.}
    \vspace{0.8em}
  \end{minipage}}
  \caption{Planned cross-benchmark view of accuracy and output compactness. The
  figure should make clear that NanoMem-GRPO improves accuracy while reducing
  the evidence passed to the answer model.}
  \label{fig:results-accuracy-compactness}
\end{figure}

\section{Compactness and Practical Efficiency}
\label{sec:results-compactness}

Compactness is not a cosmetic metric in this thesis. The memory system is
designed to be an evidence provider, so the downstream answer model should see
the minimum useful evidence state rather than a large bag of raw sessions.
However, token counts must be interpreted carefully. For NanoMem, tokens refer
to the structured evidence or memory context passed downstream. For
full-context baselines, the table columns report output-side quantities and do
not capture the full input history cost. The correct claim is therefore about
returned evidence compactness, not total system cost or latency.

The strongest compactness result is the change from NanoMem-Vanilla to
NanoMem-GRPO. On LoCoMo, mean output length drops from 264 to 107 tokens while
accuracy increases from 85.0\% to 88.6\%. On LongMemEval, output length drops
from 353 to 119 tokens while accuracy increases from 78.7\% to 83.8\%. On
ChainMem, output length drops from 287 to 234 tokens while accuracy increases
from 43.0\% to 51.5\%. This pattern suggests that training changed the
Synthesizer's evidence construction behavior, not merely the final answer
quality.

\begin{table}[H]
  \centering
  \caption{NanoMem-Vanilla versus NanoMem-GRPO compactness.}
  \label{tab:results-compactness}
  \scriptsize
  \resizebox{\textwidth}{!}{%
  \begin{tabular}{p{0.20\textwidth}rrrrrr}
    \hline\hline
    Benchmark & Vanilla acc. & Vanilla tok. & GRPO acc. & GRPO tok. & Acc. gain & Token reduction \\
    \hline
    LoCoMo & 85.0 & 264 & 88.6 & 107 & +3.6 & 59.5\% \\
    LongMemEval & 78.7 & 353 & 83.8 & 119 & +5.1 & 66.3\% \\
    ChainMem & 43.0 & 287 & 51.5 & 234 & +8.5 & 18.5\% \\
    \hline\hline
  \end{tabular}
  }
\end{table}

Compact evidence also has a risk. If the Synthesizer compresses too
aggressively, it can omit source nuance, uncertainty, or a qualifying event
needed by the answerer. The error analysis in
Section~\ref{sec:results-errors} therefore treats evidence bloat and evidence
under-specification as two sides of the same control problem.

\section{Architecture Ablation}
\label{sec:results-architecture}

The architecture ablation tests whether the gains come from the specific
structure of NanoMem rather than from simply placing another LLM before the
answer model. The planned ablations remove three mechanisms: text
normalization, the Temporal Evidence Pool, and multi-round retrieval. Each
component has a different predicted failure mode. Text normalization should
mainly affect temporal grounding. The pool should affect cross-round state
carry-over. Single-round retrieval should hurt hidden-dependency questions
more than directly addressable questions.

\begin{table}[H]
  \centering
  \caption{Architecture ablation summary. Accuracy is percentage; average
  rounds are measured on the NanoMem loop.}
  \label{tab:results-architecture-ablation}
  \scriptsize
  \resizebox{\textwidth}{!}{%
  \begin{tabular}{p{0.28\textwidth}rrrrr}
    \hline\hline
    Variant & LoCoMo-MH & LoCoMo-Temp & LME-TR & ChainMem & Avg. rounds \\
    \hline
    Full NanoMem-Vanilla & 84.0 & 76.6 & 75.8 & 43.0 & 1.25 \\
    w/o text normalization & \textcolor{red}{TODO} & \textcolor{red}{TODO} & \textcolor{red}{TODO} & \textcolor{red}{TODO} & \textcolor{red}{TODO} \\
    w/o Temporal Evidence Pool & \textcolor{red}{TODO} & \textcolor{red}{TODO} & \textcolor{red}{TODO} & \textcolor{red}{TODO} & \textcolor{red}{TODO} \\
    single-round retrieval & \textcolor{red}{TODO} & \textcolor{red}{TODO} & \textcolor{red}{TODO} & \textcolor{red}{TODO} & 1.00 \\
    \hline\hline
  \end{tabular}
  }
\end{table}

The available verified row already shows that NanoMem-Vanilla is a strong
framework before GRPO. The remaining rows must be filled from the synchronized
ablation artifact before final submission. The qualitative interpretation to
check is non-monotonic: a component need not help every task equally. For
example, single-round retrieval may be adequate for directly addressable
questions but should be weaker when the first event only reveals the next
retrieval cue.

\begin{figure}[H]
  \centering
  \fbox{\begin{minipage}{0.92\textwidth}
    \centering
    \vspace{0.8em}
    \textcolor{red}{TODO: Replace this placeholder with a heatmap of accuracy
    drop by ablation and task after the final ablation table is verified.}
    \vspace{0.8em}
  \end{minipage}}
  \caption{Planned component-impact heatmap. The expected pattern is that text
  normalization, evidence state, and multi-round search affect different task
  families.}
  \label{fig:results-architecture-heatmap}
\end{figure}

\section{Reward Diagnostics}
\label{sec:results-reward}

The reward analysis is the part of the Results chapter most tightly coupled to
the unresolved Methods marker in Section~\ref{sec:methods-reward}.
\textcolor{red}{NOTSURE: The Methods chapter currently records an unresolved
fork between the intended format/faithfulness/outcome/length reward and the
executable reward subset. Until the final training run and reward code are
synchronized, this section should be read as a diagnostic comparison of
Synthesizer behavior, not as a final causal claim about every named reward
component.}

With that caveat, Table~\ref{tab:results-reward-ablation} reports the planned
reward diagnostic values. The full reward setting has strong accuracy on the
four target tasks, high verdict accuracy, a low average number of rounds, and
short outputs. Removing the format and length terms mainly increases output
length, from about 110 to 359 tokens, while accuracy drops moderately. Removing
the verdict-related signal produces a different pattern: some temporal
categories can improve, but multi-hop and ChainMem drop, verdict accuracy falls
from 84.0\% to 72.9\%, and average rounds increase. This is consistent with
the thesis claim that stop/continue behavior is a learned part of evidence
construction.

\begin{table}[H]
  \centering
  \caption{Reward diagnostic ablation. Values should be finalized only after
  the reward implementation and reported Methods objective are reconciled.}
  \label{tab:results-reward-ablation}
  \scriptsize
  \resizebox{\textwidth}{!}{%
  \begin{tabular}{p{0.27\textwidth}rrrrrrr}
    \hline\hline
    Variant & LoCoMo-MH & LoCoMo-Temp & LME-TR & ChainMem & Verdict acc. & Avg. rounds & Avg. tok. \\
    \hline
    Full reward & 89.0 & 87.5 & 86.4 & 51.5 & 84.0 & 1.10 & 110 \\
    w/o format and length terms & 87.6 & 85.4 & 84.1 & 48.0 & 84.9 & 1.10 & 359 \\
    w/o verdict-related signal & 80.5 & 89.1 & 90.2 & 41.5 & 72.9 & 1.25 & 361 \\
    \hline\hline
  \end{tabular}
  }
\end{table}

The important interpretation is behavioral. Format and length constraints
prevent the Synthesizer from satisfying outcome reward by returning bloated
evidence. Verdict supervision helps the system distinguish between evidence
that is merely relevant and evidence that is sufficient. Without it, the loop
can stop too early on hidden-dependency questions or continue searching after
enough evidence has already been found.

\begin{figure}[H]
  \centering
  \fbox{\begin{minipage}{0.92\textwidth}
    \centering
    \vspace{0.8em}
    \textcolor{red}{TODO: Replace this placeholder with the planned reward
    ablation figure showing task accuracy, verdict accuracy, and output tokens
    as three panels.}
    \vspace{0.8em}
  \end{minipage}}
  \caption{Planned visualization of reward diagnostic behavior. The figure
  should separate the compactness failure mode from the stop/continue failure
  mode.}
  \label{fig:results-reward}
\end{figure}

\section{Qualitative Case Studies}
\label{sec:results-cases}

Aggregate scores do not show how evidence addressability is solved. The
following case-study structure should be filled with real trace examples from
the evaluation artifacts before final submission. The first case should show a
temporal bridge success: the first round finds the source event, the
Synthesizer records an event time and an \texttt{insufficient} verdict, and the
next round uses that time as a retrieval cue for the answer-bearing session.
The second case should compare NanoMem against a baseline that returns
irrelevant or overly broad context. The third case should be a NanoMem failure.

\begin{table}[H]
  \centering
  \caption{Template for a temporal bridge trace. The final thesis should
  replace this template with a real evaluation trace.}
  \label{tab:results-case-bridge}
  \scriptsize
  \resizebox{\textwidth}{!}{%
  \begin{tabular}{p{0.10\textwidth}p{0.22\textwidth}p{0.25\textwidth}p{0.16\textwidth}p{0.14\textwidth}}
    \hline\hline
    Round & Retrieval cue & Retrieved evidence summary & Event time & Verdict \\
    \hline
    1 & \textcolor{red}{TODO: source-event cue} & \textcolor{red}{TODO: source event from trace} & \textcolor{red}{TODO} & insufficient \\
    2 & \textcolor{red}{TODO: derived destination cue} & \textcolor{red}{TODO: destination evidence from trace} & \textcolor{red}{TODO} & sufficient \\
    \hline\hline
  \end{tabular}
  }
\end{table}

\textcolor{red}{TODO: Select the final case studies from real LoCoMo,
LongMemEval, or ChainMem traces. Do not fabricate query text, session IDs,
event times, verdicts, or answers. The case study must report the actual query,
key retrieved sessions, Synthesizer events, final answer, gold answer, and
diagnosis.}

\section{Error Analysis}
\label{sec:results-errors}

The observed and expected failures fall into five categories. First, retrieval
misses occur when the Planner or Retriever never returns the gold evidence
session. No Synthesizer can recover evidence it never sees. Second, temporal
grounding errors occur when a relative expression is normalized incorrectly, an
event time is inferred with the wrong span, or session time is confused with
event time. Third, premature sufficiency occurs when the Synthesizer emits
\texttt{sufficient} before all required evidence is present. Fourth,
over-search and evidence bloat occur when the system continues retrieving
after enough evidence has been collected or returns redundant evidence. Fifth,
downstream answerer errors occur when the memory evidence is adequate but the
answer model still produces the wrong final answer.

\begin{table}[H]
  \centering
  \caption{Error taxonomy for NanoMem results.}
  \label{tab:results-error-taxonomy}
  \scriptsize
  \begin{tabular}{p{0.20\textwidth}p{0.25\textwidth}p{0.22\textwidth}p{0.21\textwidth}}
    \hline\hline
    Error type & Observable signal & Likely source & Result implication \\
    \hline
    Retrieval miss & Gold session absent from selected sessions & Planner cue or retriever recall failure & Evidence synthesis is bounded by recall \\
    Temporal grounding error & Wrong date span or event ordering & Normalization or Synthesizer event-time inference & Later retrieval follows the wrong address \\
    Premature sufficiency & Verdict is sufficient before gold coverage & Synthesizer stop decision & Loop stops before bridge evidence is found \\
    Over-search or bloat & Extra rounds or long redundant evidence & Weak stop/length control & Higher cost and possible answer noise \\
    Answerer error & Evidence appears correct but answer is wrong & Downstream Qwen3.5-9B answer model & Memory and answering failures must be separated \\
    \hline\hline
  \end{tabular}
\end{table}

\textcolor{red}{TODO: Add counts or representative examples for each error
type after manually auditing the selected traces. If reliable counts are not
available, keep the analysis qualitative and state that it is based on
representative cases.}

\section{Summary}
\label{sec:results-summary}

The results support three main conclusions. First, NanoMem-GRPO is competitive
with or stronger than full-context, write-time memory, search-time retrieval,
and temporal memory baselines on the evaluated long-term memory benchmarks.
Second, the gains are largest on multi-hop, temporal reasoning, and
hidden-dependency bridge questions, which are the settings targeted by the
evidence addressability gap. Third, the trained Synthesizer improves both
accuracy and compactness: it returns shorter evidence than NanoMem-Vanilla
while improving downstream answer accuracy.

The chapter also leaves clear boundaries for the Discussion. The reward
analysis depends on synchronizing the final Methods description with the
executable reward code. The token metric measures returned evidence
compactness rather than full system cost. Case studies and error counts must be
finalized from real traces. These limitations do not remove the central result,
but they define what should be claimed carefully in the final thesis.
